\chapter{系统实现}\label{chap:impl}

本章简述将前述算法、协商协议与可视化设计落地为可运行原型时的工程组织。整套系统按职责划分为三层：数据与模型层把档案语料蒸馏为可训练张量并产出嵌入与子图模式；后端服务层把候选评分与多智能体协商封装为对外接口；前端界面层把六视图组装为单页应用。下文依次说明三层分工、与大语言模型服务的对接，以及数据与缓存策略。

\section{三层职责分工}\label{sec:impl-arch}

数据与模型层承担前期重活：从档案语料读入人年级记录，统一清洗个体、家户与事件编码；把人物、家户、社区、八旗的属性张量化，把九类亲缘与归属关系组织为异构图；按队列年构造时间因果子图，供异构图变换器与子图链路预测在其上训练，最终产出训练快照与子图模式词典。该流水线基于成熟的深度图学习框架实现\citep{paszke2019pytorch,fey2019pyg}。

后端服务层把训练产物以服务形式对外提供：对候选对返回初评分与剖面摘要，编排六轮多智能体协商，把每轮论据片段流式逐段回送前端。前端界面层把服务层的输出组织为统一的探索界面：六视图借助状态容器持有候选对与协商进度，借助事件总线转发提示路由消息。这样的分层让算法、接口与界面三类改动互不牵扯。

\section{与大语言模型服务的对接}\label{sec:impl-llm}

智能体的语言推理由外部大语言模型服务承担。后端把每一轮的角色提示、上下文与候选对档案打包为对话请求，借助流式通信通道把模型回复逐字符转发给前端，使气泡"边写边显"，避免长回复的等待。

为兼顾离线场景，后端在启动时检测大模型服务凭证：若凭证缺失或外网不可达，调用链路自动切换为本地占位智能体。占位智能体每轮返回确定性分数与套话化论据，使六视图的拓扑、动画与交互逻辑仍可演示，仅文本退化为示例口吻。

\section{数据与缓存策略}\label{sec:impl-cache}

档案数据按队列年构造时间因果子图。同一目标年的子图可能被反复读取，因此本研究为之设计磁盘缓存。缓存键由四要素组成：目标年、图模式版本、消融标记，以及当前丢弃对集合的哈希。

消融标记是必要字段：消融与未消融两条运行只改写母系亲缘边的内容，婚姻边的丢弃对集合完全相同，哈希值在两类运行下一致；若不包含消融标记，第二次运行会静默命中第一次的缓存而读到错误内容。引入消融标记后两类运行各自维护独立条目；模式或边内容变化时可经清空开关使旧缓存整体失效。